% ============================================================
% CUSTOM TIKZ MACROS FOR SPLIT ELLIPSE AND CIRCLE
% ============================================================
% This file contains reusable TikZ commands for drawing ellipses
% and circles with split styling (solid/dotted segments)
% ============================================================

% Required packages (ensure these are loaded in main document):
% \usetikzlibrary{calc}
% \usepackage{ifthen}

% ============================================================
% HELPER: \GetPointXYcm
% Get point coordinates in cm WITHOUT outer transforms
% (works correctly even with scale/rotate on tikzpicture)
%
% Usage: \GetPointXYcm{(x,y)}{\outX}{\outY}
%   Outputs \outX, \outY as decimal cm values (globally defined)
% ============================================================
\makeatletter
\newcommand{\GetPointXYcm}[3]{%
  \begin{scope}[reset cm]%
    \path #1;%
    \global\edef#2{\the\pgf@x}%
    \global\edef#3{\the\pgf@y}%
  \end{scope}%
  \pgfmathsetmacro#2{#2/1cm}%
  \pgfmathsetmacro#3{#3/1cm}%
}
\makeatother

% ============================================================
% COMMAND 1: \drawSplitEllipse (WITH ROTATION SUPPORT)
% ============================================================
% Usage:
% \drawSplitEllipse{centerX}{centerY}{rotationAngle}{semiMajor}{semiMinor}{intersectD}{intersectE}{style1}{style2}
%
% Parameters:
%   #1, #2: Ellipse center coordinates (x, y)
%   #3:     Rotation angle in degrees
%   #4:     Semi-major axis (a)
%   #5:     Semi-minor axis (b)
%   #6, #7: Two intersection points (D, E)
%   #8:     Style for first part (e.g., "Solid" or "Dot")
%   #9:     Style for second part (e.g., "Dot" or "Solid")
% ============================================================

\newcommand{\drawSplitEllipse}[9]{
    % Save the ellipse parameters
    \def\centerX{#1}
    \def\centerY{#2}
    \def\rotAngle{#3}
    \def\semiA{#4}
    \def\semiB{#5}
    
    % Save intersection points
    \coordinate (InterD) at #6;
    \coordinate (InterE) at #7;
    
    % -------------------------------------------------------
    % STEP 1: Extract x,y coordinates from points (UNTRANSFORMED)
    % -------------------------------------------------------
    \GetPointXYcm{#6}{\dX}{\dY}
    \GetPointXYcm{#7}{\eX}{\eY}
    
    % -------------------------------------------------------
    % STEP 2: Convert intersection points to local coordinates
    % Translate to center, then rotate by -rotAngle
    % -------------------------------------------------------
    % For point D:
    \pgfmathsetmacro{\dRelX}{\dX - \centerX}
    \pgfmathsetmacro{\dRelY}{\dY - \centerY}
    \pgfmathsetmacro{\dLocalX}{\dRelX * cos(-\rotAngle) - \dRelY * sin(-\rotAngle)}
    \pgfmathsetmacro{\dLocalY}{\dRelX * sin(-\rotAngle) + \dRelY * cos(-\rotAngle)}
    
    % For point E:
    \pgfmathsetmacro{\eRelX}{\eX - \centerX}
    \pgfmathsetmacro{\eRelY}{\eY - \centerY}
    \pgfmathsetmacro{\eLocalX}{\eRelX * cos(-\rotAngle) - \eRelY * sin(-\rotAngle)}
    \pgfmathsetmacro{\eLocalY}{\eRelX * sin(-\rotAngle) + \eRelY * cos(-\rotAngle)}
    
    % -------------------------------------------------------
    % STEP 3: Calculate parametric angles
    % t = atan2(y/b, x/a) - CORRECTED FORMULA
    % -------------------------------------------------------
    \pgfmathsetmacro{\angD}{atan2(\dLocalY/\semiB, \dLocalX/\semiA)}
    \pgfmathsetmacro{\angE}{atan2(\eLocalY/\semiB, \eLocalX/\semiA)}
    
    % Normalize to [0,360)
    \pgfmathsetmacro{\aD}{mod(\angD+360,360)}
    \pgfmathsetmacro{\aE}{mod(\angE+360,360)}
    
    % CCW sweep from D to E in [0,360)
    \pgfmathsetmacro{\dCCW}{mod(\aE-\aD+360,360)}
    
    % Precompute arc end angles (avoids PGF expression issues inside arc)
    \pgfmathsetmacro{\aDplus}{\aD+360}
    \pgfmathsetmacro{\aEplus}{\aE+360}
    \pgfmathsetmacro{\aEminus}{\aE-360}
    
    % -------------------------------------------------------
    % STEP 4: Determine styles based on input
    % -------------------------------------------------------
    \ifthenelse{\equal{#8}{Solid}}{%
        \tikzset{styleOne/.style={thick, solid}}%
    }{%
        \tikzset{styleOne/.style={thick, dotted}}%
    }
    
    \ifthenelse{\equal{#9}{Solid}}{%
        \tikzset{styleTwo/.style={thick, solid}}%
    }{%
        \tikzset{styleTwo/.style={thick, dotted}}%
    }
    
    % -------------------------------------------------------
    % STEP 5: Draw the two ellipse arcs with proper direction
    % -------------------------------------------------------
    \begin{scope}[shift={(\centerX,\centerY)}, rotate=\rotAngle]
        % If CCW sweep <= 180: arc1 is the short CCW arc D->E
        % else: arc1 is the short CW arc D->E (by using end<start)
        \pgfmathparse{\dCCW <= 180 ? 1 : 0}
        \ifnum\pgfmathresult=1
            % Arc 1 (short, CCW): D -> E
            \draw[styleOne] 
                (\aD:\semiA cm and \semiB cm) 
                arc (\aD:\aE:\semiA cm and \semiB cm);
            
            % Arc 2 (remaining, CCW): E -> D+360
            \draw[styleTwo] 
                (\aE:\semiA cm and \semiB cm) 
                arc (\aE:\aDplus:\semiA cm and \semiB cm);
        \else
            % Arc 1 (short, CW): D -> E (end = aE-360 < aD)
            \draw[styleOne] 
                (\aD:\semiA cm and \semiB cm) 
                arc (\aD:\aEminus:\semiA cm and \semiB cm);
            
            % Arc 2 (remaining, CCW): (aE-360) -> aD
            \draw[styleTwo] 
                (\aEminus:\semiA cm and \semiB cm) 
                arc (\aEminus:\aD:\semiA cm and \semiB cm);
        \fi
    \end{scope}
}

% ============================================================
% COMMAND 2: \drawSplitEllipseSimple (NO ROTATION)
% ============================================================
% Simplified version for ellipses with rotation angle = 0
%
% Usage:
% \drawSplitEllipseSimple{centerX}{centerY}{semiMajor}{semiMinor}{intersectD}{intersectE}{style1}{style2}
%
% Parameters:
%   #1, #2: Ellipse center coordinates (x, y)
%   #3, #4: Semi-major (a), Semi-minor (b)
%   #5, #6: Two intersection points (D, E)
%   #7, #8: Styles ("Solid" or "Dot")
% ============================================================

\newcommand{\drawSplitEllipseSimple}[8]{
    % Save parameters
    \def\centerX{#1}
    \def\centerY{#2}
    \def\semiA{#3}
    \def\semiB{#4}
    
    % Extract x,y coordinates from points (UNTRANSFORMED)
    \GetPointXYcm{#5}{\dX}{\dY}
    \GetPointXYcm{#6}{\eX}{\eY}
    
    % Calculate angles for intersection points
    % Relative to ellipse center (no rotation needed)
    \pgfmathsetmacro{\dRelX}{\dX - \centerX}
    \pgfmathsetmacro{\dRelY}{\dY - \centerY}
    \pgfmathsetmacro{\eRelX}{\eX - \centerX}
    \pgfmathsetmacro{\eRelY}{\eY - \centerY}
    
    % Correct param angles for ellipse
    \pgfmathsetmacro{\angD}{atan2(\dRelY/\semiB, \dRelX/\semiA)}
    \pgfmathsetmacro{\angE}{atan2(\eRelY/\semiB, \eRelX/\semiA)}
    
    % Normalize to [0,360)
    \pgfmathsetmacro{\aD}{mod(\angD+360,360)}
    \pgfmathsetmacro{\aE}{mod(\angE+360,360)}
    
    % CCW sweep from D to E in [0,360)
    \pgfmathsetmacro{\dCCW}{mod(\aE-\aD+360,360)}
    
    % Precompute arc end angles (avoids PGF expression issues inside arc)
    \pgfmathsetmacro{\aDplus}{\aD+360}
    \pgfmathsetmacro{\aEplus}{\aE+360}
    \pgfmathsetmacro{\aEminus}{\aE-360}
    
    % Styles
    \ifthenelse{\equal{#7}{Solid}}{%
        \tikzset{styleOne/.style={thick, solid}}%
    }{%
        \tikzset{styleOne/.style={thick, dotted}}%
    }
    
    \ifthenelse{\equal{#8}{Solid}}{%
        \tikzset{styleTwo/.style={thick, solid}}%
    }{%
        \tikzset{styleTwo/.style={thick, dotted}}%
    }
    
    \begin{scope}[shift={(\centerX,\centerY)}]
        % If CCW sweep <= 180: arc1 is the short CCW arc D->E
        % else: arc1 is the short CW arc D->E (by using end<start)
        \pgfmathparse{\dCCW <= 180 ? 1 : 0}
        \ifnum\pgfmathresult=1
            % Arc 1 (short, CCW): D -> E
            \draw[styleOne] 
                (\aD:\semiA cm and \semiB cm) 
                arc (\aD:\aE:\semiA cm and \semiB cm);
            
            % Arc 2 (remaining, CCW): E -> D+360
            \draw[styleTwo] 
                (\aE:\semiA cm and \semiB cm) 
                arc (\aE:\aDplus:\semiA cm and \semiB cm);
        \else
            % Arc 1 (short, CW): D -> E (end = aE-360 < aD)
            \draw[styleOne] 
                (\aD:\semiA cm and \semiB cm) 
                arc (\aD:\aEminus:\semiA cm and \semiB cm);
            
            % Arc 2 (remaining, CCW): (aE-360) -> aD
            \draw[styleTwo] 
                (\aEminus:\semiA cm and \semiB cm) 
                arc (\aEminus:\aD:\semiA cm and \semiB cm);
        \fi
    \end{scope}
}

% ============================================================
% COMMAND 3: \drawSplitCircle
% ============================================================
% Draw a circle split into two arcs with different styles.
% Circle is rotation-invariant, so no rotation handling needed.
%
% Usage:
% \drawSplitCircle{centerX}{centerY}{radius}{intersectD}{intersectE}{style1}{style2}
%
% Parameters:
%   #1, #2: Circle center coordinates (x, y)
%   #3:     Radius
%   #4, #5: Two intersection points (D, E) on the circle
%   #6:     Style for short arc D->E  ("Solid" or "Dot")
%   #7:     Style for remaining arc E->D  ("Solid" or "Dot")
% ============================================================

\newcommand{\drawSplitCircle}[7]{%
    % -------------------------------------------------------
    % STEP 1: Determine styles based on input
    % -------------------------------------------------------
    \ifthenelse{\equal{#6}{Solid}}{%
        \tikzset{styleOne/.style={thick, solid}}%
    }{%
        \tikzset{styleOne/.style={thick, dotted}}%
    }
    \ifthenelse{\equal{#7}{Solid}}{%
        \tikzset{styleTwo/.style={thick, solid}}%
    }{%
        \tikzset{styleTwo/.style={thick, dotted}}%
    }%
    
    % -------------------------------------------------------
    % STEP 2: Extract x,y coordinates from points D and E (UNTRANSFORMED)
    % -------------------------------------------------------
    \GetPointXYcm{#4}{\dX}{\dY}
    \GetPointXYcm{#5}{\eX}{\eY}
    
    % -------------------------------------------------------
    % STEP 3: Translate to circle's local frame (center = 0,0)
    % -------------------------------------------------------
    \pgfmathsetmacro{\xD}{\dX - (#1)}
    \pgfmathsetmacro{\yD}{\dY - (#2)}
    \pgfmathsetmacro{\xE}{\eX - (#1)}
    \pgfmathsetmacro{\yE}{\eY - (#2)}
    
    % -------------------------------------------------------
    % STEP 4: Compute parametric angles using atan2(y, x)
    % -------------------------------------------------------
    \pgfmathsetmacro{\angD}{atan2(\yD, \xD)}
    \pgfmathsetmacro{\angE}{atan2(\yE, \xE)}
    
    % Normalize to [0,360)
    \pgfmathsetmacro{\aD}{mod(\angD+360,360)}
    \pgfmathsetmacro{\aE}{mod(\angE+360,360)}
    
    % CCW sweep from D to E in [0,360)
    \pgfmathsetmacro{\dCCW}{mod(\aE-\aD+360,360)}
    
    % Precompute arc end angles (avoids PGF expression issues inside arc)
    \pgfmathsetmacro{\aDplus}{\aD+360}
    \pgfmathsetmacro{\aEplus}{\aE+360}
    
    % -------------------------------------------------------
    % STEP 5: Draw the two arcs; short arc gets styleOne
    % -------------------------------------------------------
    \begin{scope}[shift={(#1,#2)}]
        \ifdim \dCCW pt < 180.001pt
            % Short arc is CCW: D -> E
            \draw[styleOne]
                (\aD:#3) arc[start angle=\aD, end angle=\aE, radius=#3];
            % Remaining arc: E -> D+360
            \draw[styleTwo]
                (\aE:#3) arc[start angle=\aE, end angle=\aDplus, radius=#3];
        \else
            % Short arc is CW: equivalent to CCW from E -> D
            \draw[styleOne]
                (\aE:#3) arc[start angle=\aE, end angle=\aD, radius=#3];
            % Remaining arc: D -> E+360
            \draw[styleTwo]
                (\aD:#3) arc[start angle=\aD, end angle=\aEplus, radius=#3];
        \fi
    \end{scope}%
}

% ============================================================
% USAGE EXAMPLES (commented out):
% ============================================================
% 
% Example 1: Ellipse with rotation
% \drawSplitEllipse{0}{0}{0}{6}{4.472}{(-4.196,-3.196)}{(2.911,3.911)}{Solid}{Dot}
% 
% Example 2: Swap solid/dotted
% \drawSplitEllipse{0}{0}{0}{6}{4.472}{(-4.196,-3.196)}{(2.911,3.911)}{Dot}{Solid}
% 
% Example 3: Simplified version (no rotation)
% \drawSplitEllipseSimple{0}{0}{6}{4.472}{(-4.196,-3.196)}{(2.911,3.911)}{Solid}{Dot}
% 
% Example 4: Split circle
% \drawSplitCircle{0}{0}{3}{(3,0)}{(0,3)}{Solid}{Dot}
% 
% ============================================================
